% Generated draft: requires mathematical and source review.
\documentclass[11pt]{article}
\usepackage[margin=25mm]{geometry}
\usepackage{fontspec}
\setmainfont{Latin Modern Roman}
\usepackage{amsmath,amssymb,mathtools,unicode-math}
\setmathfont{Latin Modern Math}
\usepackage{enumitem,needspace,xurl,hyperref}
\setlength{\parindent}{0pt}
\setlength{\parskip}{5pt}
\setlength{\emergencystretch}{3em}
\newcommand{\N}{\mathbb{N}}
\newcommand{\Z}{\mathbb{Z}}
\newcommand{\Q}{\mathbb{Q}}
\newcommand{\R}{\mathbb{R}}
\newcommand{\C}{\mathbb{C}}
\newcommand{\F}{\mathbb{F}}
\newcommand{\A}{\mathbb{A}}
\newcommand{\PP}{\mathbb{P}}
\newcommand{\E}{\mathbb{E}}
\newcommand{\eps}{\varepsilon}
\newcommand{\dd}{\,\mathrm{d}}
\begin{document}
\section*{UnsolvedMath ID: 20002004}
\textbf{Source problem code:} catalogue entry
\subsection{Definitions and mathematical statement}
Steve Zelditch's A.22 is a collection of statistical-algebraic questions. Its background places Gaussian measures on degree-\(N\) real or complex polynomials, emphasizing projectively induced \(O(m+1)\)- and \(U(m+1)\)-invariant ensembles. It distinguishes real random counts from complex counts that may be deterministic. The principal requests are:
\begin{enumerate}
\item For random real plane algebraic curves, determine the most probable number of connected components and the expectation.
\item For random degree-\(N\) spherical harmonics, determine the most probable number of nodal domains, meaning connected components of the complement of the zero set, and its expectation.
\item In \(m\) variables, choose \(m\) random polynomials with at most a fixed \(f\) monomials, spectra consisting of \(f\) lattice points in \(\Z_{\ge0}^m\cap N\Sigma\), where \(\Sigma=\{x_i\ge0:\sum_i x_i\le1\}\), then choose coefficients from the indicated invariant ensemble. Determine the expected simultaneous real-zero count. Also determine its behavior for spectra in dilated prescribed Newton polytopes \(N\Delta_j\), with \(f\) fixed, and for the Kac alternative using normal monomial coefficients without projective monomial normalization.
\item Count and study the distribution of critical points \(\nabla s(z)=0\) of holomorphic sections \(s\in H^0(M,L)\). These depend on the Hermitian metric \(h\) or connection \(\nabla\), usually the metric connection. The source associates these with string-theoretic ``supersymmetric vacua'' and attributes the counting/statistical request to Mike Douglas.
\end{enumerate}
The ambient real domain, precise spectrum law, independence conventions, boundary-coordinate zeros, and covariance normalization require specification before these become unique numerical questions.

The source attributes a fewnomial bound to Khovanski, but its extracted expression \texttt{2m2f (f-1)/2(m+1) f} has lost exponent placement; no exact inequality is reconstructed from that typography. Its description as an enormous overestimate is a historical assessment. It also reports an exact fewnomial density formula with Shiffman, while saying asymptotics were then unavailable and Kac real-zero counts decrease. The readable kernel and constant definitions are
\[
 \Pi_N|_S(x,y)=\sum_{\beta\in S}\binom{N}{\beta}x^\beta y^\beta,
 \qquad Q=\int_{\R^m}|\xi|e^{-\langle\xi,\xi\rangle}\dd\xi,
\]
where \(\binom N\beta\) is the projective multinomial coefficient and \(\mathcal S(N,f)\) is the set of spectra. The displayed density averages over \(S\) and involves a square root of a determinant of \(\nabla_x\nabla_y\log\Pi_N|_S\) on the diagonal, but its prefactor and placement of the kernel powers are corrupted. No exact density identity is asserted here. Shared versus independent spectra for the equations also require checking.

The aid's exact partial result is conditional on a fixed binomial system on the real algebraic torus:
\[
 c_i x^{a_i}+d_i x^{b_i}=0\quad(1\le i\le m),\qquad
 x\in(\R^*)^m,
\]
where \(a_i,b_i\in\Z^m\), the integer matrix \(B\) with rows \(a_i-b_i\) is invertible over \(\R\), and coefficients are nonzero independent random variables with mutually independent unbiased signs. Set \(r=\operatorname{rank}_{\F_2}(B\bmod2)\). If \(Z\) counts all real-torus roots, the reported law is
\[
 \Pr(Z=2^{m-r})=2^{r-m},\qquad
 \Pr(Z=0)=1-2^{r-m},\qquad \E Z=1.
\]
Every root is simple. If \(Z_+\) counts roots in \((\R_{>0})^m\), the aid gives \(\E Z_+=2^{-m}\). These formulas do not count roots on coordinate hyperplanes. They depend on sign assumptions, not a particular continuous distribution of coefficient magnitudes.

If all exponents of this fixed system are multiplied by an integer \(D\ge1\), odd \(D\) preserves the law, while even \(D\) gives
\[
 \Pr(Z_D=2^m)=2^{-m},\qquad\Pr(Z_D=0)=1-2^{-m},\qquad\E Z_D=1.
\]
This deterministic exponent dilation is distinct from resampling spectra in \(N\Delta_j\). The binomial theorem alone does not answer arbitrary \(f\)-fewnomial, curve-topology, or holomorphic-section questions.

For the usual rotationally invariant Gaussian degree-\(N\) harmonic on \(S^2\), the aid additionally reports from Nazarov--Sodin concentration that every choice of a mode \(q_N\) of its nodal-domain count satisfies \(q_N\sim c_{\mathrm{NS}}N^2\), with the corresponding positive constant in that normalization. No precise concentration estimate or citation is supplied. Mere convergence in probability by itself need not control modes of integer-valued variables, so the quantitative input in that reported derivation needs review.
\subsection{Short English statement}
The record asks several questions about random real curves, nodal domains, sparse polynomial zeros, and metric-dependent critical points of holomorphic sections. The aid gives an exact real-torus binomial law, whose distribution changes under even exponent dilation although its mean stays one. It also reports a spherical-harmonic mode asymptotic; neither claim resolves all of the original statistical questions.
\subsection{Sources}
\begin{itemize}
\item Local UnsolvedMath registry, permanent ID 20002004, exact source code catalogue entry; canonical \texttt{aim-geometry-notes.json}, \texttt{notes[341]}, tag \texttt{section}; AIM Workshop Problem Lists, Geometry, \emph{Amoebas and tropical geometry}, A.22, contributed by Steve Zelditch. Created and updated 14 August 2026.
\item AIM source \url{https://aimath.org/WWN/amoebas/amoebas.pdf}. The supplied verification identifies a collection of related questions rather than one assertion. The PDF was not accessed.
\item Original text names Khovanski for a fewnomial bound, Shiffman with the contributor for an exact density, and Mike Douglas for the holomorphic-section counting request. It mentions Picard--Lefschetz, Milnor--Orlik, and Arnold as background on critical points of holomorphic functions. No complete bibliographic entries are supplied.
\item Supplied aid names Nazarov--Sodin concentration without a precise publication or estimate. Attempt 1: \texttt{partial\_result}, problem status at run \texttt{partially\_solved}, rigor label \texttt{proved}, novelty confidence \texttt{low}; classification \texttt{PARTIAL-PROGRESS}, registry status \texttt{partially\_solved}. Candidate: exact binomial distribution synthesis and parity-sensitive dilation corollary. All status claims are dated supplied assessments.
\item Ulam AI / UnsolvedMath Contributors; CC BY 4.0. Machine-generated research aid, not peer reviewed; independently verify before citation.
\end{itemize}
\end{document}
